\chapter{实验方法}\label{chap:method}

\section{仿真环境与任务设置}

本文使用动态立方体抓取环境作为实验平台。机器人为固定基座机械臂，目标为边长$0.05\,\mathrm{m}$的立方体。中等难度真值环境由\path{config/env_gt_medium.yaml}定义，仿真步长为$0.01\,\mathrm{s}$，目标区域中心为$(0.20,0.00,0.60)\,\mathrm{m}$，目标区域半径为$0.10\,\mathrm{m}$，每个回合最长2400步。立方体初始位置在$x\in[0.15,0.35]$、$y\in[-0.35,0.35]$、$z=0.25$范围内采样，线速度在平面内采样，角速度绕固定$z$轴采样，角速度大小为$15$至$30\,\mathrm{deg/s}$。

环境观测包含立方体位置、姿态、线速度、角速度，末端执行器位置、姿态，夹爪开合状态，以及机器人关节位置。专家策略可以使用仿真真值观测；含噪实验则通过感知包装器在观测路径上加入高斯噪声、漂移、离群点、丢帧和延迟等扰动。成功判定由环境在闭环执行中给出，要求系统完成接近、抓取、保持和运输相关阶段。

\section{专家策略与数据采集}

专家策略采用有限状态机组织任务阶段，包括approach、dwell、close、hold和transport等阶段。approach阶段将末端执行器移动到目标相对抓取位姿附近；dwell阶段在目标运动中维持相对对齐；close阶段执行夹爪闭合；hold和transport阶段确认稳定抓取并将目标送往目标区域。专家内部结合任务几何构造参考位姿，并由MPC求解器在任务空间和关节空间之间执行跟踪。

该专家不是人类遥操作专家，而是“算法即专家”的合成示教源。这样做有两个直接优点：其一，仿真真值可以提供目标速度和姿态，便于生成相位清晰的示教；其二，批量采集可以覆盖大量初始位置、速度和角速度条件，降低人类示教成本。中等任务原始数据集位于\path{dataset/expert_medium}，元数据记录总回合数为12500，总帧数为2934103。记录的动作标签包括目标位置、目标姿态和夹爪命令，观测标签包括立方体、末端执行器、夹爪和关节状态。

\section{数据编译与表示}

扩散策略不直接读取完整原始回合，而是由编译器把成功回合转换成定长训练窗口。配置\path{config/dp_medium.yaml}中设定历史观测窗口$H_o=16$，动作预测窗口$H_a=32$，未来条件窗口为50，步长仍为$0.01\,\mathrm{s}$。编译过程只保留成功回合，并采用grasp-only设置聚焦抓取相关片段。最终编译数据位于\path{dataset/compiled_diffusion_medium}，包含11909个回合、387931个训练窗口和43036个验证窗口。

表~\ref{tab:dataset-summary}总结了中等任务数据规模。

\begin{table}[!htbp]
    \bicaption{中等任务数据集规模。}{Dataset scale for the medium task.}
    \label{tab:dataset-summary}
    \centering
    \small
    \begin{tabular}{lrr}
        \hline
        数据项 & 数量 & 来源 \\
        \hline
        原始专家回合 & 12500 & \path{dataset/expert_medium/meta.json} \\
        原始帧数 & 2934103 & \path{dataset/expert_medium/meta.json} \\
        编译成功回合 & 11909 & \path{dataset/compiled_diffusion_medium/manifest.json} \\
        训练窗口 & 387931 & \path{dataset/compiled_diffusion_medium/manifest.json} \\
        验证窗口 & 43036 & \path{dataset/compiled_diffusion_medium/manifest.json} \\
        \hline
    \end{tabular}
\end{table}

动态抓取中的主要学习难点之一是世界坐标动作分布随目标平移和自旋快速变化。如果直接学习世界坐标中的末端参考，同一种抓取意图会对应大量数值差异较大的动作标签。本文采用目标中心化表示：将末端执行器相对位置、相对姿态、立方体局部角速度和夹爪状态组合为历史观测特征；将未来目标运动外推表示为相对位移和相对姿态变化；将专家末端参考同样转换为相对目标坐标系下的动作。

编译后的每个历史观测维度为11，包括相对末端位置3维、相对末端四元数4维、目标局部角速度3维和夹爪状态1维。未来条件每步7维，包括相对位移3维和相对四元数4维。动作每步8维，包括相对目标位置3维、相对目标四元数4维和夹爪命令1维。

设$t$时刻的历史观测特征为$o_t$，未来目标条件为$c_t$，专家动作标签为$a_t$，则本文使用的学习样本可表示为
\begin{equation}
    (o_{t-H_o+1:t}, c_{t:t+H_c-1}) \rightarrow a_{t:t+H_a-1},
\end{equation}
其中$H_o=16$，$H_c=50$，$H_a=32$。这种形式使策略在生成当前动作时同时利用过去观测和短时未来目标外推信息。

\section{扩散策略设计}

本文策略采用条件扩散策略框架。扩散模型在训练时学习从带噪动作块恢复专家动作块，在推理时从随机噪声开始逐步去噪，得到未来动作序列。模型结构由三部分组成：历史编码器、未来条件编码器和一维U-Net去噪网络。历史编码器使用因果一维卷积提取过去16步观测；未来条件编码器处理50步目标运动条件；U-Net在扩散时间步、历史嵌入和未来条件嵌入的调制下预测动作噪声。除连续动作去噪头外，模型还包含close事件预测和阶段分类分支，用于表达抓取任务中的离散相位信息。

运行时，策略使用10步DDIM反向采样生成动作块，并通过时间集成平滑重叠动作块。策略输出的是任务空间参考，而不是最终关节命令；下层MPC继续负责把参考转换为可执行的机器人运动。这一“学习规划器加确定性执行器”的划分使学习模型专注于动态拦截和抓取对齐，同时保留MPC对执行约束和运输阶段的处理能力。

\section{训练与评估设置}

中等任务策略由\path{config/dp_medium.yaml}统一描述。模型规模约为12.6M参数，训练目标包括动作噪声预测损失、close事件二分类损失和阶段分类损失。训练使用AdamW优化器、余弦学习率调度、指数滑动平均权重和自动混合精度。训练完成后的运行检查点为\path{checkpoints/diffusion_medium/best.pt}。在本文实验中，该检查点与\path{config/env_gt_medium.yaml}共同构成扩散策略的主要评估对象。

本文比较两类策略。第一类为MPC专家，由\path{config/policy_expert.yaml}配置，代表直接使用当前观测和手工阶段逻辑的模型控制基线。第二类为扩散模仿策略，由\path{config/dp_medium.yaml}配置，运行检查点为\path{checkpoints/diffusion_medium/best.pt}。两类策略在相同环境变体上运行，每个配置使用200个评估回合和4个并行worker。

评估矩阵分为真值观测和含噪观测两部分。真值观测矩阵位于：

\begin{center}
\path{config/thesis_eval/gt_medium}
\end{center}

该矩阵包含完整中等任务以及$x$向平移、$y$向平移、平面平移、旋转、平移加旋转等分解场景。含噪矩阵位于：

\begin{center}
\path{config/thesis_eval/noisy_medium}
\end{center}

该矩阵包含轻度、中度、强度完整噪声以及中等强度下的高斯、漂移、离群点、丢帧延迟等单项扰动。本文首先报告成功率，即成功回合数除以总回合数，并使用Wilson区间给出二项成功率的95\%置信区间。为避免只用成功率描述闭环行为，评估日志还记录最终末端到立方体距离、末端跟踪参考的平均误差、跟踪误差低于阈值的比例、末端轨迹二阶差分得到的平滑度指标、成功回合步数、归一化完成速度，以及完整闭环吞吐量。吞吐量以每秒仿真步数和相对$100\,\mathrm{Hz}$控制频率的实时倍率表示，只反映当前仿真和推理实现的整体运行速度，不等同于真实硬件实时控制能力。